% GENERATED FILE. Edit canonical lessons, then run npm run generate:books.
% canonical-source-hash: fnv1a64:94c021d9fadc8388
% canonical-lessons: ES-C60-repaso-grado

\chapter{Review --- One Scale, One Habit}
\label{ch:repaso-grado}
\hlchaptermodality{229}

\begin{chapteropening}
\textbf{By the end of this chapter:} \emph{I can place a thing on a scale, and recognise a shortened word.}

I can place a thing on a scale, and recognise a shortened word.
\end{chapteropening}

\section[(review)]{(review)}

\label{lesson:ES-C60-repaso-grado}



\begin{quote}
The scale, least to most. (\emph{Poco}, \emph{bastante}, \emph{muy}.)
\end{quote}

\begin{grammarlens}[title={What you've built}]
Three chants. Say each, then say what it is for.

\textbf{First.} \emph{Estoy muy cansado.} — \emph{\textbf{Muy}} is the top of the scale, and it is \emph{mucho} with its ending bitten off.

\textbf{Second.} \emph{Estoy bastante cansado.} — \emph{\textbf{Bastante}} is the middle, and it means \emph{sufficing}: \emph{bastar} plus the \emph{-ante} you could already read.

\textbf{Third.} \emph{Estoy mal.} — \emph{\textbf{Mal}} is the opposite of \emph{bien}, and it is \emph{malo} with its ending bitten off.
\end{grammarlens}

\begin{grammarlens}[title={Grammar Lens: the habit underneath}]
Two of those three came in a long form and a short one, and that is the thing worth carrying out of this chapter.

\emph{Mucho} keeps its ending in front of a noun and loses it in front of an adjective. \emph{Malo} keeps its ending standing alone and loses it in front of a masculine noun. Different words, one habit: \textbf{the ending goes when the word is leaning on the next one.}

So when you meet a short Spanish word that looks like a longer one you already know, that is usually what has happened — not a new word to learn.
\end{grammarlens}

\subsection*{Your turn}
Say these aloud:

\begin{itemize}
\raggedright
  \item the three chants, in order
  \item \emph{mucho} / \emph{muy}, then \emph{malo} / \emph{mal}
\end{itemize}

\emph{Twice through:}

\begin{tcolorbox}[breakable,colback=teal!4,colframe=teal!35!black,title={Before you move on}]
Which of the three refuses to stand in front of a noun? (\emph{\textbf{Muy}}.)

And which one refuses to stand alone as an answer? (\emph{\textbf{Muy}} again — \emph{muy bien}, never just \emph{muy}.)
\end{tcolorbox}
